% [原文 p.9]
\chapter{有理曲线与标准类}

本章考察非奇异射影簇 $X$ 的标准类 $K_X$ 与 $X$ 上有理曲线之间的关系。

§1.1 处理 $-K_X$ 丰富的情形。这样的 $X$ 称为 Fano 簇。主要结果是：$X$ 含有关于 $-K_X$ 次数较低的有理曲线。这一结果由 [Mor79] 给出，是极小模型理论的出发点之一。即便在复簇的情形，其证明也要经过正特征，这一点颇有意思。

§1.2 不假定 $-K_X$ 丰富，只假设它与某条曲线的相交数为正，从而推广上述结果。证明只在“不可忘记若干附加信息”这一点上有所不同，其余与 §1.1 的论证极其相似。

这些有理曲线存在定理至今已有若干改进，但证明的本质部分仍然是这里所介绍的论证。

本书中，§1.1–§1.2 的论证通向 §1.3 非奇异簇锥定理的几何证明。该锥定理是 [Mor82] 的结果。第 3 章将用完全不同的方法证明更一般的锥定理。

§1.4 通过构造曲面的极小模型，展示锥定理用法的一个例子。本书其余部分的主题，正是如何把这些结果推广到高维。

最后一节给出若干基本丰富性判别法的证明。

% [原文 p.10]
§1.3 的方法不足以在高维完成极小模型理论，在其余各章中不再使用。尽管如此，这些方法向我们展示了非常明快的几何图景，将成为其后更技术性的结果的指针。

本章所说明的几何思想成为进一步研究的原动力，其中许多在 [Kol96] 中有阐述。

\section{有理曲线的寻找（标准类为负的情形）}

\takeaway{本节寻找关于 $-K_X$ 次数较低的有理曲线。所用工具是形变理论与折弯引理。结论是：当 $-K_X$ 丰富（即 $X$ 为 Fano 簇）时，过 $X$ 的每一点都有关于 $-K_X$ 次数不超过 $\dim X+1$ 的有理曲线。}

本节可以当作热身来读。这里要探讨的是下面这个一般性主题。

\textbf{1.1}\quad 簇上的有理曲线，会对该簇的双有理几何产生怎样的影响？\hfill$\square$

% [Tutor Note: 1.1 是纲领性问题，本章各节都在为它给出局部答案：无有理曲线时（如一般型簇）双有理几何相对简单；反之 Fano 簇上极丰富的有理曲线恰是极小模型纲领要收缩的对象。]

首先我们将会看到，若不存在有理曲线，状况会变得非常称心。随后反过来利用这一点，说明簇 $X$ 的双有理几何的某种复杂性，正是由 $X$ 上特殊的有理曲线引起的。

最简单的例子是曲面理论中关于 $(-1)$-曲线的 Castelnuovo 定理（参看 [Har77, V.5.7]）。

\begin{theorem}{1.2}
设 $X$ 为固有的非奇异曲面。此时，存在到非奇异曲面的、非同构的固有双有理态射 $f:X\to Y$，当且仅当 $X$ 上存在自相交数 $(C^2)=-1$ 的非奇异有理曲线 $C$。这样的 $C$ 称为 $(-1)$-曲线。
\end{theorem}

该结果的一个侧面可以容易地作如下推广。

\begin{propositionp}[{[Abh56, Prop. 4]}]{1.3}
设 $X$ 为任意维数的非奇异簇，$f:Y\to X$ 为固有双有理态射，$x$ 为 $X$ 的任意点。则 $f^{-1}(x)$ 若不是一点，便被有理曲线覆盖。
\end{propositionp}

\begin{proof}
首先处理 $X$ 为曲面的情形。反复爆破（blow up）$X$ 的点可解消 $f^{-1}$ 的不定点（indeterminacy），而每个阶段都得到 $\PP^1$。因此 $f^{-1}(x)$ 是这些 $\PP^1$ 的若干并集的像。于是由 Lüroth 定理，$f^{-1}(x)$ 是有理曲线的并集。

% [原文 p.11]
一般情形下，只要知道如何解消双有理映射的不定点，也能同样证明；但实际上远弱的形式的不定点解消就已足够。另外，后面用到命题 1.3 的只是曲面情形，故高维情形仅给出证明概略。

不妨假设 $Y$ 正规。由 [Sha94, II.4.4]，$f$ 的例外集是纯余维数 1 的（参看推论 2.63）。设 $E\subset Y$ 为例外集的不可约分支。在 $E$ 的一般点 $e$ 的近旁，$(Y,E)$ 可由 $X$ 出发、沿非奇异中心作有限次爆破得到（参看后面的引理 2.45）。

因此存在有理曲线 $C\subset E$，它含 $e$ 且 $f(C)$ 为一点。有理曲线即便退化，也只会成为有限条有理曲线的并集，故对 $E$ 的任一点，都存在含该点的 $E$ 内有理曲线。
\end{proof}

% [Tutor Note: 命题 1.3 的证明只用到“例外集纯余维 1”这种很弱的不定点解消；原书强调高维时无需一般的强解消定理，这正是省略论证的关键所在。]

（Zariski 主定理 [Har77, V.5.2]）设 $X$ 为正规簇，$f:Y\to X$ 为固有双有理态射。若 $f^{-1}$ 在点 $x\in X$ 处不正则，则 $f^{-1}(x)$ 连通且维数 $\ge 1$。

利用它，由命题 1.3 立即得到下述结果。

\begin{corollary}{1.4}
设 $Z$ 为非奇异簇，$g: Z\dashrightarrow X$ 为有理映射。$Y\subset Z\times X$ 为 $g$ 的图像闭包，$p$、$q$ 分别为到 $Z$、$X$ 的投影。$S\subset Z$ 为 $g$ 的不定点全体构成的闭子集。此时，对 $q(p^{-1}S)$ 的任一点，都存在过该点、位于 $q(p^{-1}S)$ 内的有理曲线。
\end{corollary}

\begin{corollary}{1.5}
设 $X$、$Z$ 为代数簇，$Z$ 非奇异，$X$ 固有。若存在带不定点的有理映射 $g: Z\dashrightarrow X$，则 $X$ 含有有理曲线。
\end{corollary}

能应用此推论的最简单情形，是 $Z$ 为作为曲线族而得到的曲面之时。只要把情形设置得当，就能使推论 1.5 中的有理映射 $g$ 不成为态射。

\begin{lemmap}[{刚性引理（rigidity lemma）}]{1.6}
设 $Y$ 为不可约簇，固有态射 $f:Y\to Z$ 为满射。$f$ 的任一纤维连通，维数为定值 $n$。（由上述假设，任一纤维必然为纯 $n$ 维。）设 $g:Y\to X$ 为满足下述条件的态射：对 $Z$ 的某点 $z_0$，$g(f^{-1}(z_0))$ 为一点。此时，
% [原文 p.12]
对 $Z$ 的任一点 $z$，$g(f^{-1}(z))$ 都是一点。
\end{lemmap}

\begin{proof}
$f\times g: Y\to Z\times X$ 固有，故其像 $W\subset Z\times X$ 是 $Z$ 上的固有簇。由此得到到固有态射的分解
\[ f:\ Y\xrightarrow{h}W\xrightarrow{p}Z, \]
并且 $p^{-1}(z)=g(f^{-1}(z))$、$\dim p^{-1}(z_0)=0$ 成立。由纤维维数的上半连续性，可适当选取含 $z_0$ 的开集 $U\subset Z$，使 $U$ 的每点 $z$ 上的纤维 $p^{-1}(z)$ 全为 0 维。因此 $h$ 在 $p^{-1}(U)$ 上的纤维维数为 $n$，而 $h$ 在 $W$ 上的各纤维维数至少为 $n$。由假设 $h^{-1}(w)\subset f^{-1}(p(w))$ 且 $\dim f^{-1}(p(w))=n$。故对所有 $w\in W$ 有 $\dim h^{-1}(w)=n$。另一方面 $f^{-1}(z)=h^{-1}(p^{-1}(z))$，所以 $\dim p^{-1}(z)=0$ 对所有 $z\in Z$ 成立。$f^{-1}(z)$ 连通，故如所断言，$p^{-1}(z)=h(f^{-1}(z))$ 是一点。
\end{proof}

% [Tutor Note: 引理 1.6 中“纤维连通”与“纤维维数恒为 $n$”缺一不可：若维数不恒定，上半连续性只给出开集上的信息，推不出每个纤维都映为一点。]

\begin{corollary*}[{折弯引理（bend and break）I}]{1.7}
设给定固有代数簇 $X$、固有非奇异曲线 $C$、一点 $p\in C$，以及非常值态射 $g_0: C\to X$。又设（不必固有的）非奇异连通曲线 $D$ 与点 $0\in D$，连同态射 $G: C\times D\to X$ 满足下述条件。
\begin{enumerate}[label=(\arabic*),leftmargin=2.4em,itemsep=1pt,topsep=2pt]
\item $G|_{C\times\{0\}}=g_0$；
\item $G(\{p\}\times D)=g_0(p)$；
\item 对一般的 $t\in D$，$G|_{C\times\{t\}}$ 与 $g_0$ 不同。
\end{enumerate}
此时，存在（可能为常值的）态射 $g_1: C\to X$，以及 $X$ 上有理曲线的线性组合 $Z=\sum a_iZ_i$（$a_i>0$，$Z_i\subset X$），满足下面两个条件。
\begin{enumerate}[label=(\arabic*),leftmargin=2.4em,itemsep=1pt,topsep=2pt]
\item 作为 1-循环，$(g_0)_*(C)\equiv (g_1)_*(C)+Z$。（1-循环与数值等价参看定义 1.16。）
\item $g_0(p)\in\bigcup_i Z_i$。
\end{enumerate}
特别地，$X$ 含有过 $g_0(p)$ 的有理曲线。
\end{corollary*}

本书没有认真处理代数等价，故上面的推论中改用数值等价；但作为代数等价也成立，证明照原样可行。

引用这一结果时，下面的说法虽不准确，却很好地表达了其气氛。

% [原文 p.13]
“若 $g_0$ 固定一点而移动，则它退化为有理曲线之和与另一态射 $g_1: C\to X$ 之和。”

\begin{proof}[推论 1.7 的证明]
将 $D$ 紧化为固有曲线 $\bar D$，并把态射 $G$ 延拓为有理映射 $\bar G: C\times\bar D\dashrightarrow X$。下面证明 $\bar G$ 的不定点存在于 $\{p\}\times\bar D$ 之中。为此取 $p$ 在 $C$ 中的开邻域 $U$，把引理 1.6 用于投影 $f: U\times\bar D\to U$。纤维 $f^{-1}(p)$ 被 $g$ 映为一点，故每条纤维都被映为一点。于是对一般的 $t\in D$ 也有 $G(z,t)=g_0(z)$，矛盾。这说明 $\bar G$ 在 $\{p\}\times\bar D$ 的某处有不定点。

记 $\bar G$ 的图像闭包的正规化为 $S$。设 $\pi: S\to C\times\bar D$ 与 $G_S: S\to X$ 为投影，$h: S\to C\times\bar D\to\bar D$ 为 $\pi$ 与到第二分量投影的合成。

如上所述，在 $C\times\bar D$ 的某点 $(p,d)$ 上方 $\pi$ 不是同构。除子 $h^*([d])$ 有分解 $h^*([d])=C'+E$，其中 $C\cong C'\subset S$ 是 $C\times\{d\}$ 的双有理变换，$E$ 是 $\pi$-例外的。令 $g_1: C\to X$ 为 $G_S$ 在 $C'$ 上的限制，并记 $Z=(G_S)_*E$。

$g_0$ 可与 $G_S|_{h^*(0)}$ 等同视之。因此由 $h^*([0])\equiv h^*([d])$ 得 $(g_0)_*(C)\equiv (g_1)_*(C)+Z$（参看定义 1.16）。由命题 1.3，$E$ 是有理曲线的并集，故 $Z$ 亦然。利用推论 1.4 即知存在过 $g_0(p)$ 的有理曲线。
\end{proof}

值得注意的是，上述结果中代数簇这一假设不可或缺。

\begin{example}[{[Bla56]}]{1.8}
设 $E$ 为椭圆曲线，$M$ 为次数 $\ge 2$ 的线丛，它由截面 $s$、$t$ 生成。

在直和 $V=M\oplus M$ 中，截面 $(s,t)$、$(\sqrt{-1}\,s,-\sqrt{-1}\,t)$、$(t,-s)$、$(\sqrt{-1}\,t,\sqrt{-1}\,s)$ 处处 $\R$-线性无关。因此它们生成 $E$ 上的一个格丛（lattice bundle）$L$。令 $X=V/L$，$C$ 为 $V/L$ 的零截面。此时由于向量丛 $V$ 为正，$C$ 可以固定一点而移动。然而 $V/L$ 不含有理曲线。

结论：在 $C\to X$ 的态射的形变族中，固定 $C$ 的一点的那些，除平凡情形外都不能紧化。
\end{example}

% [Tutor Note: 例 1.8 说明定理 1.10 对非代数的紧复流形失效：$V/L$ 是复环面，$V$ 的正性使曲线可以固定一点而移动，它却没有有理曲线。故证明必须经由正特征而不能停在复解析范畴。]

% [原文 p.14]
当推论 1.7 中的曲线 $C$ 为有理曲线时，可取 $g_1:=g_0$，故推论 1.7 是平凡的。推论 1.7 的下一个变体说明如何得到有理曲线的非平凡退化（degeneration）。

\begin{lemma*}[{折弯引理 II}]{1.9}
设给定射影簇 $X$ 与非常值态射 $g_0:\PP^1\to X$。又设（不必固有的）非奇异连通曲线 $D$、点 $0_D\in D$ 和态射 $G:\PP^1\times D\to X$ 满足下述条件。
\begin{enumerate}[label=(\arabic*),leftmargin=2.4em,itemsep=1pt,topsep=2pt]
\item $G|_{\PP^1\times\{0_D\}}=g_0$；
\item $G(\{0\}\times D)=g_0(0)$，$G(\{\infty\}\times D)=g_0(\infty)$；
\item $G(\PP^1\times D)$ 是曲面。
\end{enumerate}
此时，$(g_0)_*(\PP^1)$ 在 $X$ 上数值等价于可约曲线或重数曲线（即形如 $aC$、系数 $a>1$ 的曲线）（参看定义 1.16）。
\end{lemma*}

与推论 1.7 一样，注意本引理的结论作为代数等价也成立。

\begin{proof}
设 $\bar D$ 为 $D$ 的非奇异紧化，$q: S\to\bar D$ 为以 $\PP^1\times D$ 为开集含于其中、且 $q$ 延拓第二投影 $\PP^1\times D\to D$ 的 $\PP^1$-丛（例如直积），$\bar G: S\dashrightarrow X$ 为 $G$ 的延拓。设 $\tilde r:\tilde S\to S$ 为若干次单点爆破，使由 $\bar G$ 决定的有理映射 $\tilde G:\tilde S\dashrightarrow X$ 成为态射。对 $\tilde r$ 中爆破的次数作归纳来证明引理。

先处理 $\bar G$ 已是态射的情形。设 $H$ 为 $X$ 上的丰富除子，$C_0,C_\infty\subset S$ 为延拓 $\{0\}\times D$、$\{\infty\}\times D$ 的 $q$ 的截面。则 $((\bar G^*H)^2)>0$，且 $(C_0\cdot\bar G^*H)=0=(C_\infty\cdot\bar G^*H)$。由 Hodge 指数定理，$(C_0^2)<0$，$(C_\infty^2)<0$。另一方面 $S$ 是 $\bar D$ 上的 $\PP^1$-丛，故 $S$ 的 Néron–Severi 群（参看定义 1.16）的秩为 2（例 1.23(1)）。有效曲线 $C_0$、$C_\infty$ 与 $\bar G^*H$ 正交，故存在 $a,b>0$ 使 $aC_0\equiv bC_\infty$。于是 $(C_0^2)=(b/a)(C_0\cdot C_\infty)\ge 0$，矛盾。

% [Tutor Note: 此处省略一步：$C_0$、$C_\infty$ 都与 $\bar G^*H$ 正交，而 $(\bar G^*H)^2>0$，故二者落在其正交补中；正交补在秩 2 的 Néron–Severi 群中是一维的，于是两曲线成比例。]

设 $\bar G$ 不是态射，取不定点解消所需点 $P\in q^{-1}(y)$ 处的单点爆破
\[ \tilde S\xrightarrow{\tilde r}S'\xrightarrow{r}S. \]
设 $F_1$ 为 $r$ 的例外曲线。则 $(q\circ r)^*(y)=F_1+F_2$，其中 $F_1$、$F_2$ 是交于一点 $Q=F_1\cap F_2$ 的 $(-1)$-曲线。

% [原文 p.15]
设 $\bar G': S'\dashrightarrow X$ 为自然确定的有理映射。

下面证明 $\bar G'$ 在 $F_2$ 的每一点处正则。由于 $(g_0)_*(\PP^1)\equiv\tilde G_*((q\circ\tilde r)^*(y))$，为证明引理，注意可以假定右端作为 1-循环只有一个成分且系数为 1。若 $G$ 在 $q^{-1}(y)$ 中还有别的不定点 $P\ne P'\in q^{-1}(y)$，则 $\tilde G_*\operatorname{red}(\tilde r^{-1}(P))$ 与 $\tilde G_*\operatorname{red}(\tilde r^{-1}(P'))$ 都含有曲线（推论 1.4），且 1-循环
\[ \tilde G_*((q\circ\tilde r)^*(y))-\tilde G_*\operatorname{red}(\tilde r^{-1}(P))-\tilde G_*\operatorname{red}(\tilde r^{-1}(P')) \]
为有效（各系数 $\ge 0$），矛盾。下面证明 $S'\dashrightarrow X$ 在 $Q$ 处有定义。$\operatorname{red}((r')^{-1}(Q))$ 的所有不可约分支在 $(q\circ\tilde r)^*(y)$ 中的重数都 $\ge 2$，因此必须被 $\tilde G$ 收缩。断言至此得证。

设 $S'\to S''$ 为 $(-1)$-曲线 $F_2$ 的收缩，则有理映射 $S''\dashrightarrow X$ 用少一次爆破即可解消不定点。由归纳法，引理得证。
\end{proof}

陈述并证明关于有理曲线存在性的第一个主要结果，准备已经就绪。该定理至今已被 [MM86]、[Miy87]、[Kaw91]、[K+92, 第 9 章] 等改进，但证明的本质部分未变。

\begin{theoremp}[{[Mor79]}]{1.10}
设 $X$ 为非奇异射影簇，$-K_X$ 丰富。则 $X$ 含有有理曲线。事实上，对 $X$ 的每一点 $x$，都存在过它的有理曲线 $D$ 满足
\[ 0<-(D\cdot K_X)\le\dim X+1. \]
\end{theoremp}

\begin{proof}
证明由若干步骤组成。

\textbf{步骤 1}\ 我们想应用推论 1.7。因此需要找出可形变的态射 $f: C\to X$。任取曲线 $C$。不过，若想找到过 $X$ 的任意给定点 $x$ 的有理曲线，就应选取 $C$ 过 $x$，并取 $C$ 上一点 $0$ 使 $f(0)=x$。

\textbf{步骤 2}\ 使用形变理论中下述非平凡的结果。关于形变理论的许多解说书（例如 [Kol96, II.1] 等）中都有这一解说。

（维数估计公式 [Mor79]）从亏格（genus）为 $g(C)$ 的射影曲线 $C$ 到非奇异簇 $X$ 的态射 $f$ 的形变空间，其维数满足下述估计。
% [原文 p.16]
\begin{equation*}
\begin{aligned}
\text{维数}&\ge h^0(C,f^*T_X)-h^1(C,f^*T_X)\\
&=-(C\cdot f^*K_X)+(1-g(C))\cdot\dim X.
\end{aligned}
\end{equation*}

这一结果通过下述考察来使用。

若形变空间的维数为 $m$，则存在 $m$ 维不可约仿射簇 $Z\ni 0$ 与态射 $F: C\times Z\to X$，使得 $F|_{C\times\{0\}}=f$，并且对任意 $z_1\ne z_2\in Z$ 有 $F|_{C\times\{z_1\}}\ne F|_{C\times\{z_2\}}$。

固定点 $0$ 在 $f$ 下的像 $f(0)$ 至多是 $\dim X$ 个条件，故把 $0$ 映到 $x$ 的态射 $f: C\to X$ 的形变空间维数有如下下界估计
\begin{equation*}
\begin{aligned}
\text{维数}&\ge h^0(C,f^*T_X)-h^1(C,f^*T_X)-\dim X\\
&=-(C\cdot f^*K_X)-g(C)\cdot\dim X.
\end{aligned}
\end{equation*}
因此，若 $-(C\cdot f^*K_X)-g(C)\cdot\dim X>0$，则存在态射 $f: C\to X$ 的非平凡单参数族，且固定 $f(0)$。由引理 1.6，得到过 $x$ 的 $X$ 的有理曲线。证明的这一部分对 Kähler 流形也有效，但如例 1.8 所示，对任意的紧非奇异复流形并不成立，这一点请注意。

% [Tutor Note: 形变空间维数估计本质是“取负号的相交数”：$-(C\cdot f^*K_X)$ 越大，形变空间超出固定基点所需的 $\dim X$ 个条件就越多，曲线因而被迫移动。]

\textbf{步骤 3}\ 说明如何达到条件 $-(C\cdot f^*K_X)-g(C)\cdot\dim X>0$。为此必须使 $-(C\cdot f^*K_X)$ 足够大。分三种情形讨论。
\begin{enumerate}[label=(\arabic*),leftmargin=2.4em,itemsep=1pt,topsep=2pt]
\item $g(C)=0$。若 $-(C\cdot f^*K_X)>0$，则 $C$ 在 $X$ 上移动。但已经知道 $X$ 上有过 $x$ 的有理曲线。
\item $g(C)=1$。若 $-(C\cdot f^*K_X)>0$，就把 $f$ 与 $C$ 的 $n$ 倍自同态合成，其中 $n$ 为整数。于是
\[ -(C\cdot(f\circ n)^*K_X)-\dim X=-n^2(C\cdot f^*K_X)-\dim X. \]
因此 $f\circ n: C\to X$ 固定一点而移动。
\item $g(C)\ge 2$。亏格不小于 2 的曲线没有次数大于 1 的自同态。于是考虑另一条曲线 $h: C'\to C$ 并形变 $f\circ h$。此时的问题在于：$(C'\cdot(f\circ h)^*K_X)=\deg h\cdot(C\cdot f^*K_X)$ 虽能变大，但 $C'$ 的亏格也随之改变。
\end{enumerate}
例如，即便想让 $C$ 的 $m$ 重非分歧覆盖 $C'$ 移动，形变空间

% [原文 p.17]
的维数估计式
\[ \text{维数}\ge m[-(C\cdot f^*K_X)-g(C)\cdot\dim X]+(m-1)\dim X \]
只是得到了保证。无论 $m$ 取多大，即便在 $-(C\cdot f^*K_X)>0$ 时，这个下界也未必变大。

\textbf{步骤 4}\ 如上所见，在 $g(C)>1$ 的情形，$C$ 没有高次自同态，因而这一尝试不奏效。不过，确有 $C$ 实际上具有高次自同态的情形，那就是正特征的情形，此时 Frobenius 态射正是所需的自同态。下面看如何从给定的情形转移到正特征 $p>0$ 的情形。

\textbf{步骤 5}\ 取非奇异簇 $X\subset\PP^n$ 与带有到 $X$ 的态射的曲线 $C\subset\PP^m$。首先假设 $C$ 与 $X$ 都由整系数方程定义：
\[ h_1(x_0,\dots,x_n)=\cdots=h_r(x_0,\dots,x_n)=0\ \text{定义}\ X, \]
\[ c_1(y_0,\dots,y_m)=\cdots=c_s(y_0,\dots,y_m)=0\ \text{定义}\ C. \]
设 $\mathbb{F}_p$ 为含 $p$ 个元素的有限域，$\bar{\mathbb{F}}_p$ 为其代数闭包。此时上述方程组 $h_i$、$c_j$ 分别在射影空间 $\bar{\mathbb{F}}_p\PP^n$ 与 $\bar{\mathbb{F}}_p\PP^m$ 中定义簇 $X_p$、$C_p$。

对几乎所有的 $p$（即除 $p$ 的有限个例外值外），这些簇非奇异且 $\dim C_p=1$。映射
\[ (x_0,\dots,x_n)\mapsto(x_0^p,\dots,x_n^p) \]
定义了 $C_p$ 的 Frobenius 自同态 $F_p$。它作为集合间的映射是单射，却被看作次数为 $p$ 的态射。

% [Tutor Note: 页图印作 $x_0,\dots,x_n$，但由 $C_p\subset\PP^m$ 的上下文看应为 $y_0,\dots,y_m$ 的坐标；英文原著作 $y$。此处照页图照录。]

另一方面，由 $\operatorname{Spec}\Z$ 上的生成平坦性（generic flatness）[Mat69, Thm. 52]，对几乎所有的 $p$，态射 $C\to\operatorname{Spec}\Z$ 在 $p$ 的逆像上平坦。因此对这些 $p$，$(C_p\cdot(f_p)^*K_{X_p})$、$g(C_p)$ 等值是一定的（参看 [Har77, III.9.10]）。这里态射
\[ f_p\circ F_p^m:\ C_p\xrightarrow{F_p^m}C_p\xrightarrow{f_p}X_p \]
的附带基点条件的形变空间，其维数估计式为
\[ \text{维数}\ge -p^m(C_p\cdot(f_p)^*K_{X_p})-g(C_p)\cdot\dim X. \]

% [原文 p.18]
由于 $(C_p\cdot(f_p)^*K_{X_p})$ 是与 $p$ 无关的负常数，上述下界 $-p^m(C_p\cdot(f_p)^*K_{X_p})-g(C_p)\cdot\dim X$ 当 $m\gg 1$ 时为正（对几乎所有的 $p$，取 $m=1$ 就已足够）。于是与步骤 2 一样，对几乎所有的 $p$ 都能构造 $X_p$ 上的有理曲线 $A_p$。

\textbf{步骤 6}\ 下面处理更一般的情形：$h_j$（$X$ 在 $\PP^n$ 中的定义方程）、$c_j$（$C$ 在 $\PP^m$ 中的定义方程）以及 $b_j$（态射图像在 $\PP^n\times\PP^m$ 中的定义方程）等的系数未必是整数。这些系数生成基域的有限生成部分环 $R\supset\Z$。注意：若 $P$ 是 $R$ 的任一极大理想，则 $R/P$ 是有限域（无限域不可能作为 $\Z$ 上的环有限生成）。因此适当选取 $p$ 后，$R/P$ 同构于某个 $\mathbb{F}_{p^k}$，即含 $p^k$ 个元素的有限域。此时 Frobenius 态射由把 $\mathbb{F}_{p^k}\PP^n$ 的齐次坐标 $(x_0,\dots,x_n)$ 分别取 $p^k$ 次幂给出。其余论证同上进行，在 $\operatorname{Spec}R$ 的某稠密 Zariski 开集的任一几何闭点 $P$ 上，都构造出有理曲线 $A_p$。

\textbf{步骤 7}\ 现设 $-K_X$ 丰富，且 $X$ 由 $|-mK_X|$ 嵌入，其中 $m$ 为适当正整数。这里想把 $A_p$ 换成为满足 $-(K_{X_p}\cdot B_p)\le\dim X+1$ 的另一条有理曲线 $B_p$。

为此注意：若 $-(K_{X_p}\cdot A_p)>\dim X+1$，则态射 $A_p\to X_p$ 在固定两点 $q,q'$ 之下至少有二维的形变族。然而 $\PP^1$ 固定两点的自同构只有一维，故 $X_p$ 中 $A_p$ 的像也会移动。于是用引理 1.9 可知 $A_p$ 分解为更低次的有理曲线之和。因此只要 $-(A_p\cdot K_X)>\dim X+1$，就存在次数更低的有理曲线。

\textbf{步骤 8}\ 这最后一步中，必须由几乎所有的 $p$ 都存在有界次数的有理曲线 $B_p$ 这一事实，导出特征 0 的整系数簇 $X$ 上存在有理曲线。（一般情形用 $P\in\operatorname{Spec}R$ 同样论证。）
\end{proof}

% [边界: 原理 1.11 的陈述首句起于 p.18、止于 p.19，按 BRIEF §4 规则 1 在本块内写完。]
\begin{principlep}{1.11}
设整系数的齐次多项式系对无限多个素数 $p$ 在 $\bar{\mathbb{F}}_p$ 上有非平凡解（系数非整数时，改为对 $\operatorname{Spec}R$ 的某 Zariski 稠密子集中所有的几何点 $p$ 有非平凡解）。此时，该方程组在任意代数闭域上都有非平凡解。
\end{principlep}
